\newglossaryentry{g_Objective(OKR)}{
	name={Objective/s},
	description={
		Der Begriff \textbf{Objective} bezieht sich in diesem Kontext auf die OKR-Logik. Im Gegensatz zum militärischen Gebrauch im Englischen ist es hier ein wünschenswerter Zustand, der nicht zwingend eine eindeutige Zielerreichung voraussetzt. Dieser Zustand kann kontinuierlich angepasst und überprüft werden. Ein Beispiel: \textit{Mein Ziel ist, dass Deutschland ein demokratisches Land ist.} Was genau unter Demokratie verstanden wird, kann zu Beginn klar definiert sein, sich aber im Laufe der Zeit ändern und mit neuen Erkenntnissen abgeglichen werden.
		Daher ist es bei einem \textit{Objective} besonders wichtig, die \textit{Key Results} und \textit{Maßnahmen} präzise zu beschreiben. Diese definieren, welche Indikatoren den Erfolg messen oder den Fortschritt hin zum \textit{Objective} bestimmen. Die \textit{Maßnahmen} legen konkret fest, welche Schritte unternommen werden, um das \textit{Objective} letztendlich zu erreichen.
	}
}

\newglossaryentry{g_Rubric}{
    name={Rubric},
    description={
        A \textit{rubric} is a coherent collection of principles, goals, or theses that are integrated under a \gls{p_CDF} (Coherent Decision Framework). 
        It defines a topic by organizing and aligning its elements so that they mutually reinforce each other and avoid contradictions. 
        The development of a rubric includes defining the conceptual boundaries of the topic and ensuring coherence through alignment, non-contradiction, and integration of its parts. 
        A rubric serves as a structured framework for reasoning and decision-making, which can be communicated clearly and consistently. 
        In contrast to a Strategy, a rubric does not prescribe specific actions but provides the foundational logic and coherence for them.
    }
}